\begin{figure}[ht]
\centering
\begin{tikzpicture}[
  land/.style = {circle, draw=engblack, thick, fill=englime!25,
                 inner sep=0pt, minimum size=9mm, font=\small},
]
% Konigsberg multigraph: four landmasses, seven bridges.
% N (north bank), S (south bank), I (Kneiphof island), E (east land).
\node[land] (N) at (2, 3)   {$N$};
\node[land] (S) at (2, -3)  {$S$};
\node[land] (I) at (0, 0)   {$I$};
\node[land] (E) at (4.4, 0) {$E$};

% Two bridges each between I and N, I and S; one each elsewhere.
\draw[thick, draw=engblue] (I) to[bend left=25]  (N);
\draw[thick, draw=engblue] (I) to[bend right=25] (N);
\draw[thick, draw=engblue] (I) to[bend left=25]  (S);
\draw[thick, draw=engblue] (I) to[bend right=25] (S);
\draw[thick, draw=engblue] (I) -- (E);
\draw[thick, draw=engblue] (N) -- (E);
\draw[thick, draw=engblue] (S) -- (E);

% Degree annotations.
\node[font=\scriptsize, anchor=south] at (2, 3.55)  {$\deg = 3$};
\node[font=\scriptsize, anchor=north] at (2, -3.55) {$\deg = 3$};
\node[font=\scriptsize, anchor=east]  at (-0.55, 0) {$\deg = 5$};
\node[font=\scriptsize, anchor=west]  at (4.95, 0)  {$\deg = 3$};
\end{tikzpicture}
\caption{The seven bridges of K\"{o}nigsberg as a multigraph: four
landmasses (north bank $N$, south bank $S$, Kneiphof island $I$,
east land $E$) joined by seven bridges. All four vertices have odd
degree. Euler's 1736 argument shows that a walk crossing every edge
exactly once (an \engterm{Euler trail}) exists only if the number
of odd-degree vertices is $0$ or $2$; here it is $4$, so no such
walk exists. The negative result, that the townsfolk's Sunday walk
was impossible, founded graph theory. The engineer reads the
odd-degree count, not the map, to settle the question.}
\label{fig:vol02:ch17:konigsberg}
\end{figure}
